\documentclass[tikz,border=5pt]{standalone}
\usetikzlibrary{positioning, calc}

\definecolor{myblue}{RGB}{204, 204, 255}
\definecolor{myyellow}{RGB}{255, 255, 204}

\nopagecolor
\pagecolor{black!90} 


\begin{document}
\begin{tikzpicture}[
    cell/.style={
        draw=white, 
        thick,
        minimum height=1.3cm,
        align=center,
        font=\scriptsize\ttfamily,
        inner sep=4pt,
        text=white,
        outer sep=0pt,
        node distance=0pt
    },
    % 3:1 width ratio for 48-bit vs 16-bit
    bluecell/.style={cell, fill=none, minimum width=4.5cm},
    yellowcell/.style={cell, fill=none, minimum width=2.2cm}, % 2.2cm gives comfortable padding for "Element Index"
    dimline/.style={draw=white, thin},
    dimlabel/.style={text=white, font=\scriptsize\sffamily, inner sep=3pt},
    bitnum/.style={font=\scriptsize\ttfamily, text=white}
]

% 48-bit Chunk Index
\node[bluecell] (chunk) {Chunk Index \\ (48 bits)};

% 16-bit Element Index
\node[yellowcell, right=of chunk] (element) {Element Index \\ (16 bits)};

\node[bitnum, below=2pt of chunk.south west]   {63};
\node[bitnum, below=2pt of element.south west, xshift=6pt] {15};
\node[bitnum, below=2pt of element.south east] {0};
\node[bitnum, below=2pt of chunk.south east, xshift=-6pt] {16};



% 64-bit bracket dimension line across the entire top
%\draw[dimline] (chunk.north west) -- ++(0, 4mm) -- node[dimlabel] {64 bits} ([yshift=4mm]element.north east) -- (element.north east);

%\draw[dimline] (chunk.north west) -- ++(0, 4mm) -- node[above, text=white, font=\scriptsize\sffamily] (lbl) {64 bits Total} ([yshift=4mm]element.north east) -- (element.north east);
    
\node[text=white, font=\scriptsize\sffamily] (L) at ([yshift=4mm]$(chunk.north west)!0.5!(element.north east)$) {64 bits};
\draw[dimline] (chunk.north west) |- (L) (element.north east) |- (L);
       

\end{tikzpicture}
\end{document}
